\section*{Chapter \ref{ch:g3:typing-with-ten-fingers} --- Typing with Ten Fingers}

\begin{solution}{exo:g3:typing-with-ten-fingers:1}
A: left little finger. S: left ring finger. E: left middle finger. R and T:
left index finger. J: right index finger. O: right ring finger. P: right little
finger. (On an AZERTY keyboard the A is on the top row, still under the left
little finger; on a QWERTZ keyboard these eight letters do not move.)
\end{solution}

\begin{solution}{exo:g3:typing-with-ten-fingers:2}
Capital A: the right \key{Shift}, with the right little finger. Capital L: the
left \key{Shift}, with the left little finger.
\end{solution}

\begin{solution}{exo:g3:typing-with-ten-fingers:3}
They let the index fingers find the \omterm{def:g3:typing-with-ten-fingers:home-row}{home row} by touch, without looking at the
keys; the other fingers then fall into place.
\end{solution}

\begin{solution}{exo:g3:typing-with-ten-fingers:4}
$150 \div 5 = 30$ words in 2 minutes, so $30 \div 2 = 15$ \omterm{def:g3:typing-with-ten-fingers:wpm}{words per minute}.
\end{solution}

\begin{solution}{exo:g3:typing-with-ten-fingers:5}
$30 \times 3 = 90$ words, and $90 \times 5 = 450$ characters.
\end{solution}

\begin{solution}{exo:g3:typing-with-ten-fingers:6}
For example: looking down at the keys; typing with two fingers only; leaving
the \omterm{def:g3:typing-with-ten-fingers:home-row}{home row} and searching for it again; making mistakes that must be rubbed
out.
\end{solution}

\begin{solution}{exo:g3:typing-with-ten-fingers:7}
QWERTY and AZERTY: A, Q, Z, W and M are not in the same place (A and Q are
swapped, Z and W are swapped, M moves to the \omterm{def:g3:typing-with-ten-fingers:home-row}{home row}). QWERTY and QWERTZ: Y
and Z are swapped.
\end{solution}

\begin{solution}{exo:g3:typing-with-ten-fingers:8}
H: right index finger (second column); o: right ring finger; m: right index
finger; e: left middle finger.
\end{solution}

\begin{solution}{exo:g3:typing-with-ten-fingers:9}
Mia: she types accurately, and speed will come. Tom's 6 mistakes must each be
rubbed out and typed again, and mistakes made while learning become habits.
\end{solution}

\begin{solution}{exo:g3:typing-with-ten-fingers:10}
15 hours are $15 \times 60 = 900$ minutes, and $900 \div 10 = 90$ days.
\end{solution}

\begin{solution}{exo:g3:typing-with-ten-fingers:11}
The index fingers are the strongest and the easiest to move, so they take the
two middle columns. The little fingers are at the edges of the hands, so they
take the keys at the edges of the keyboard.
\end{solution}
